% ============================================================
%  part5/ch32-agency.tex  —  Chapter 32: Simulated Agency
% ============================================================

\chapter{Simulated Agency}
\label{ch:agency}


\begin{ChapterIntro}
An agent is often imagined as something fundamentally distinct from the environment around it.

This chapter takes a different approach.

Instead of treating agency as a special substance, it treats agents as structured regions within a larger field. An agent differs from its surroundings not because it escapes the laws governing them, but because it participates in a particular kind of recursive loop: it observes, reconstructs, and acts, and its actions alter the world it subsequently observes.

The challenge is to formalise this loop without introducing new ontological categories. If agency is real, it should emerge naturally from the same admissibility structures that govern every other aspect of the framework.
\end{ChapterIntro}

\section{The agent as admissible subsystem}

\begin{definition}[Agent]
  \label{def:agent}
  A \emph{simulated agent} is a tuple
  \[
    A = (U, \mathcal{A}_U, F_A, G_A, \pi_A, \alpha_A),
  \]
  where $U \subseteq \mathcal{M}$ is the agent's local domain,
  $F_A$ decomposes observations, $G_A$ reconstructs a
  world-model, $\pi_A$ projects to actionable states, and
  $\alpha_A$ maps reconstructions to field perturbations.
\end{definition}

\begin{definition}[Admissible action space]
  Let $\mathcal{E}$ be the subsheaf of compactly supported
  admissible perturbations over $U$.
  The \emph{viable action space} is $H^0(U, \mathcal{E})$.
\end{definition}

\begin{proposition}[Agent paralysis theorem]
  \label{prop:paralysis}
  \statusproven{}
  If $H^0(U, \mathcal{E}) = 0$, no nontrivial admissible
  action exists.
\end{proposition}

\begin{proof}
  $H^0(U, \mathcal{E})$ is the space of admissible local sections.
  If it contains only the zero section, $\alpha_A = 0$ for every
  admissible action: no nontrivial perturbation can be generated.
\end{proof}

\begin{InterpBox}
  The agent is not a special substance.
  An agent is a region capable of constructing a local model
  and acting through admissible perturbations.
  Agency is not added to the framework as a new ontological
  category: it is a particular class of admissible section.
\end{InterpBox}

\section{Agency requires spatial extent}

\begin{proposition}[Geometric capacity for agency]
  For an agent region $U$ to possess nontrivial action sections,
  $U$ must be large enough to support a smooth perturbation
  satisfying both internal constraints and boundary decay.
\end{proposition}

This provides a minimum geometric capacity for agency:
a point-like agent is formally stripped of its capacity to
alter the field, because no smooth perturbation can be
constructed in a region too small to smooth out its gradients
before reaching the boundary.

\section{Multi-agent interactions}

When two agent regions $U_1$ and $U_2$ overlap
($U_{12} = U_1 \cap U_2 \neq \emptyset$), their actions
generate a mismatch measured by $H^1(U_1 \cup U_2, \mathcal{E})$.

\begin{theorem}[Agent conflict as gluing obstruction]
  \label{thm:agent-conflict}
  \statusproven{}
  If $\alpha_A^{(1)}$ and $\alpha_A^{(2)}$ are admissible actions
  on $U_1$ and $U_2$, and their restrictions to $U_{12}$ disagree,
  then their combination defines a nontrivial class in
  $H^1(U_1 \cup U_2, \mathcal{E})$.
\end{theorem}

\principle{
  Agent conflict is not a new force.\\[4pt]
  It is a nonzero gluing obstruction\\[4pt]
  relaxed by RSVP free-energy descent.
}

\begin{definition}[Conflict resolution functional]
  \[
    \delta\Phi^*_{12}
    = \operatorname*{argmin}_{s \in \mathcal{F}(U_{12})}
    \left[
      \|\delta\Phi_1|_{U_{12}} + \delta\Phi_2|_{U_{12}} - s\|^2
      + \beta\,\mathcal{F}_{\mathrm{RSVP}}[s]
      + \gamma\,\mathcal{O}_{\mathrm{glue}}[s]
    \right].
  \]
\end{definition}

Conflict resolution proceeds through constrained minimisation
of the same free-energy/admissibility structure governing all
RSVP dynamics.
The existing nonlinearities (energy boundedness, entropy production,
vorticity dissipation) handle attenuation without requiring a
new physical force.

\section{The recursive agency loop}

The agency condition is recursive:
\[
  G_A(F_A(U)) \to \text{action} \to U' \to G_A(F_A(U')).
\]
Agency is a closed admissibility loop inside RSVP.
Not a separate ontology: an admissibility-preserving subsystem
whose actions remain gluable into the surrounding sheaf.

\begin{MathEcho}
  This chapter introduced:
  $A = (U, \mathcal{A}_U, F_A, G_A, \pi_A, \alpha_A)$,
  $H^0(U, \mathcal{E})$ (action space),
  agent paralysis (when $H^0 = 0$),
  and conflict as $H^1$ obstruction.
  Agency is the self-referential case of the reconstruction
  framework: an agent is a system that applies $G \circ F$
  to itself and acts on the result.
\end{MathEcho}

\section{Agency existence criterion and paralysis theorem}

\begin{theorem}[Agency existence criterion]
  \label{thm:agency-existence}
  \statusproven{}
  Agency is possible (i.e., nontrivial admissible action exists)
  if and only if $H^0(U, \mathcal{E}) \neq 0$.
\end{theorem}

\begin{proof}
  By definition, $H^0(U,\mathcal{E})$ is the space of global
  admissible sections of $\mathcal{E}$ over $U$, i.e., the
  viable action space.
  If $H^0(U,\mathcal{E}) \neq 0$, there exists a nontrivial
  section $\alpha_A \neq 0$, so nontrivial action is possible.
  If $H^0(U,\mathcal{E}) = 0$, the only section is $\alpha_A = 0$,
  so no nontrivial perturbation can be generated.
\end{proof}

\begin{theorem}[Paralyzed agent theorem]
  \label{thm:paralyzed-agent}
  \statusproven{}
  If $H^0(U,\mathcal{E}) = 0$, then every admissible action
  satisfies $\alpha_A = 0$: the agent is formally paralyzed.
\end{theorem}

\begin{proof}
  This is the contrapositive of the existence criterion:
  if $H^0(U,\mathcal{E}) = 0$, no nontrivial admissible action
  exists, so $\alpha_A = 0$ for all admissible $\alpha_A$.
\end{proof}

\begin{InterpBox}[What paralysis means physically]
  Agent paralysis is not a failure of will or intention.
  It is a topological fact about the ambient field.
  If the region $U$ is too small, too constrained, or too
  tightly embedded in a rigid field configuration, no smooth
  perturbation can be constructed that satisfies both the
  internal state requirements and the boundary admissibility
  conditions.
  Agency requires \emph{room}: sufficient geometric capacity
  to generate a nontrivial smooth perturbation that respects
  the surrounding sheaf structure.
\end{InterpBox}

% ---- Figures and citations ----------------------------------

\[
\begin{tikzcd}[column sep=3em]
  \text{Observe}
  \arrow[r]
  &
  \text{Reconstruct}
  \arrow[r,"G(F(\cdot))"]
  &
  \text{Act}
  \arrow[r,"\alpha_A"]
  &
  \text{Observe}
  \arrow[loop above]{}{}
\end{tikzcd}
\]

\[
\begin{tikzcd}
  H^0(U,\mathcal{E})
  \arrow[r,"\alpha_A"]
  &
  \mathcal{F}(U)
  \arrow[r,"\text{gluing}"]
  &
  \mathcal{F}(X)
\end{tikzcd}
\]

\begin{CitationAnchor}
  Similar recursive observe-reconstruct-act feedback structures
  appear throughout cybernetics (Ashby 1952, Wiener 1948),
  control theory (Powers 1973), and modern predictive-processing
  frameworks (Friston~\cite{friston2010}).
  The formulation here differs by requiring actions to lie in
  $H^0(U,\mathcal{E})$ --- the global sections of the admissible
  perturbation sheaf --- which is a genuinely new condition
  not present in earlier frameworks.
\end{CitationAnchor}

\begin{FurtherReading}
  Cybernetics and feedback-based agency are developed in
  Wiener's \textit{Cybernetics} (1948) and Ashby's
  \textit{Design for a Brain} (1952).
  Friston's free-energy principle for predictive processing
  appears in Friston~\cite{friston2010} and provides a
  contemporary counterpart.
  For the mathematical treatment of reconstruction as Bayesian
  inference, see Tenenbaum et al.~\cite{tenenbaum2011}.
\end{FurtherReading}
